Language-model agents read tool outputs written for programs: most of what an API returns is never visibly reused,
yet is paid for at every turn. We learn from an agent's own traces which fields of each tool output it visibly
reuses and remove the rest, protecting identifiers and letting the agent recover any hidden field with an
\texttt{expand} tool. In live runs of tracetrim 0.4 on \texttt{\SthreeModelA}, under frozen protocols, a fixed
projection with \texttt{expand} lowered the median LLM API cost per run by \SthreeEOneTrimCosti{} and
\SthreeEOneTrimCostiv{} on two agents with every implemented gate passing, though a stricter audit found billing
diagnoses less complete with trimming, so preservation of answer quality is not yet established; when a changed task
needed removed fields, adaptive learning kept them from the next run in \SthreeHFiveKept{} of
\SthreeHFiveRecoveries{} cases. Offline, with today's defaults, the learned projection removed \BaseLearnedRawi,
\BaseLearnedRawiii{} and \BaseLearnedRawv{} of held-out tool-output tokens on three of five agents, and
\BaseLearnedStaticBigMin{} to \BaseLearnedStaticBigMax{} of what a minifier and a static metadata filter fixed in
advance leave there. Earlier studies (versions 0.1 and 0.2) explain these defaults: learned from tool calls alone, as
a gateway would, the 0.1 aggressive projection missed a used field in \MissesPooledShareRealCalls{} of the held-out
structured calls it trims on the real-API agents, against \MissesPooledShareRealFull{} from complete traces; replays
located three failure mechanisms (transformed values, fields used only to reason, invented identifiers), which the
safeguards of 0.2 repaired on the conversations already observed. The method helps most where outputs carry much the
task does not use, least where the agent cites most of what it reads; it has not yet been validated on tasks held out
from development. The library, a telemetry report and an MCP proxy are released.
